\par\Needspace{6\baselineskip}\medskip
\begin{center}\hypertarget{native-170-article123}{}\pdfbookmark[2]{123.}{article-123}123.\end{center}
Bisher haben wir vorausgesetzt, dass Näherungswerte der Unbekannten $x,y$ bereits aus irgendeiner Quelle vorliegen. Sind die ungefähren Abmessungen der ganzen Bahn bekannt, etwa aus früheren Rechnungen mit anderen Beobachtungen gewonnen und nun anhand neuer Beobachtungen zu berichtigen, so lässt sich diese Bedingung ohne Schwierigkeit erfüllen, welche Bedeutung wir den Unbekannten auch geben. Bei der ersten Bestimmung einer noch völlig unbekannten Bahn, der bei Weitem schwierigsten Aufgabe, kommt es dagegen sehr darauf an, welche Unbekannten wir verwenden. Sie sind mit Bedacht so zu wählen, dass sich ihre Näherungswerte aus der Natur der Aufgabe selbst gewinnen lassen. Dies gelingt am besten, wenn die drei zur Untersuchung der Bahn verwendeten Beobachtungen eine nicht zu große heliozentrische Bewegung des Himmelskörpers umfassen. Solche Beobachtungen sind daher stets für die erste Bestimmung zu verwenden, die man dann nach Belieben anhand weiter auseinanderliegender Beobachtungen berichtigen kann. Es ist leicht einzusehen, dass die unvermeidlichen Beobachtungsfehler die Rechnung umso stärker stören, je näher die verwendeten Beobachtungen beieinanderliegen. Für die erste Bestimmung sind die Beobachtungen daher mit Bedacht auszuwählen. Man muss \emph{erstens} darauf achten, dass sie nicht zu dicht beieinanderliegen, und \emph{zweitens}, dass sie nicht zu weit auseinanderliegen. Lägen die Beobachtungen zu dicht beieinander, so ließen sich die ihnen genügenden Elemente zwar sehr leicht berechnen, doch könnte man sich auf diese Elemente selbst kaum verlassen. Die Fehler könnten sie sogar so stark entstellen, dass sie nicht einmal als Näherung brauchbar wären. Lägen die Beobachtungen zu weit auseinander, so fehlten uns die Hilfsmittel zur näherungsweisen Bestimmung der Unbekannten. Wir erhielten dann \SourceLeaf{171}nur eine sehr grobe Näherung, für die viel mehr Hypothesen nötig wären, oder eine ganz ungeeignete, bei der sich äußerst mühsame Versuche nicht vermeiden ließen. Diese Grenzen des Verfahrens richtig zu beurteilen lernt man besser durch häufige Übung als durch Vorschriften. Die unten folgenden Beispiele werden zeigen, dass sich aus Beobachtungen der Juno, die nur 22 Tage auseinanderliegen und eine heliozentrische Bewegung von $7^0 35'$ umfassen, bereits Elemente mit beträchtlicher Genauigkeit gewinnen lassen. Umgekehrt lässt sich unser Verfahren noch mit sehr gutem Erfolg auf Beobachtungen der Ceres anwenden, die 260 Tage auseinanderliegen und eine heliozentrische Bewegung von $62^0 55'$ umfassen. Mit vier Hypothesen oder vielmehr vier aufeinanderfolgenden Näherungen liefert es Elemente, die mit den Beobachtungen sehr gut übereinstimmen.
